
In this chapter, I will show some formulas omitted from Chapter \ref{ch:micro}.


\begin{comment}


The first step is to set the perturbation theory for the Hubbard-Kanamori Hamiltonian
\begin{equation}
     { H } _{HK} =  { H } _{\mathrm{at}}+ { H } _\text{hop}
\end{equation}
in the limit $U,J_H\gg t$. The combination of average one-hole occupancy per site and large $U$ drives the electronic system to a Mott insulating phase. For finite $U$, the holes can jump from one site/orbital to the next, lowering their kinetic energy. %Pron this virtual exchange process  the low energy description of this theory
In the single-orbital Hubbard model, the well-known superexchange mechanism  favors the antiferromagnetic arrangement of neighboring spin. Projecting into the single occupancy subspace, the charge degrees of freedom remain quenched and the effective Hamiltonian is a spin-1/2 Heisenberg Hamiltonian  $J \bm{\sigma}_i\cdot\bm{\sigma}_j$ with $J\sim t_{\scriptsize 1,3}^2/U$. Here we will apply the standard procedure of perturbation theory for $ { H } _{HK}$.


%put a high energy cost for double occupancy of holes

To implement the perturbation theory we will use the canonical transformation
\begin{equation}
     { H } _{HK} \to \tilde{ { H } } = e^{\, S}  \;   { H } _{HK} \, e^{-S} =  { H } _{HK} + [S, { H } _{HK}]+\frac{1}{2}[S,[S, { H } _{HK}]] + \cdots \, , \label{eq:II-1-S-transf}
\end{equation}
which eliminates perturbation in the first order, that is the terms that change $\bar{N}$ up to order $(t/U)^2$. We have used the Baker-Campbell-Hausdorff formula to expand the exponentials in the right-hand-side of equation \eqref{eq:II-1-S-transf}.

%Effective Hamiltonian in the low-energy subspace up to third order in $t/U$.
We can use the projector $\mathcal{P}^{(n)}_{i}$ to the subspace of $n$ holes at the site $i$. In this low-energy sector, we can approximate the identity by $\mathbb{1} \approx \mathcal{P}^{(0)}_{ }+\mathcal{P}^{(1)}_{ }+\mathcal{P}^{(2)}_{ }$, which differ from the sector of one hole per site at most by one hole. The hopping Hamiltonian can be written as
\begin{equation}
     { H } _{\text{hop}} = - \sum_{ij} \mathbb{1} \;  h_i^\dagger  \tilde{\mathbf{T}}_{ij} h_j^{\phantom{\dagger}}  \, \mathbb{1} \approx T_0 + T_1 + T_{-1}\; ,
\end{equation}
where I am defining $T_m$ as the term that changes the hole occupation in the hopping Hamiltonian by $m$
\begin{align}
\begin{split}
    T_0 &= \mathcal{P}^{(0)}  T \mathcal{P}^{(0)} + \mathcal{P}^{(1)}  T \mathcal{P}^{(1)} +\mathcal{P}^{(2)}  T \mathcal{P}^{(2)}  \\
    &= \mathcal{P}^{(1)}  T \mathcal{P}^{(1)} +\mathcal{P}^{(2)}  T \mathcal{P}^{(2)}  \; ,
\end{split}\\
\begin{split}
    T_1 &=  \mathcal{P}^{(1)}  T \mathcal{P}^{(0)} +\mathcal{P}^{(2)}  T \mathcal{P}^{(1)}
    \\ &=\mathcal{P}^{(2)}_{}  T \mathcal{P}^{(1)}    \; ,
    \end{split}\\
\begin{split}
    T_{-1} &= \mathcal{P}^{(0)}  T \mathcal{P}^{(1)} +\mathcal{P}^{(1)}  T \mathcal{P}^{(2)}   \; ,
     \\ &=\mathcal{P}^{(1)} T \mathcal{P}^{(2)} \; .
    \end{split}
\end{align}
Here we use that  $ \mathcal{P}^{(n)}_{i}  h_{is}^{\dagger}  \mathcal{P}^{(m)}_{i} $ is non-zero only if $n=m+1$, and $h_i\mathcal{P}^{(0)}_{i}=\mathcal{P}^{(0)}_{i}h_i^\dagger=0$.
Explicitly we have
\begin{align}
    T_0 &=   - \sum_{ij}  \mathcal{P}^{(1)}_{i}   h_i^\dagger  \tilde{\mathbf{T}}_{ij} h_j^{\phantom{\dagger}}  \mathcal{P}^{(1)}_{j} - \sum_{ij}\mathcal{P}^{(2)}_{i}   h_i^\dagger  \tilde{\mathbf{T}}_{ij} h_j^{\phantom{\dagger}}  \mathcal{P}^{(2)}_{j}    \; , \\
    T_1 &=   - \sum_{ij}\mathcal{P}^{(2)}_{i}   h_i^\dagger  \tilde{\mathbf{T}}_{ij} h_j^{\phantom{\dagger}}  \mathcal{P}^{(1)}_{j}    \; , \\
T_{-1} &= - \sum_{ij} \mathcal{P}^{(1)}_{i}   h_i^\dagger  \tilde{\mathbf{T}}_{ij} h_j^{\phantom{\dagger}} \mathcal{P}^{(2)}_{j} = T_1^{\dagger}    \; .
\end{align}
The decomposition of the hopping Hamiltonian in these different hopping operators allows us to separate the terms in the transformed Hamiltonian $\tilde{H}$ according to the contribution for the change in the occupation number.

%\subsection{Commutators}
\subsection{Canonical transformation}

For the canonical transformation, we need to calculate the commutators between the hopping operators $T_m$ and the atomic Hamiltonian $ { H } _{\mathrm{at}}$. First, we split the projector operator for two-hole states into channels with different orbital angular momentum
\begin{equation}
    \mathcal{P}^{(2)}_{i} = \mathcal{P}^{(2)}_{i,\ell=0}+\mathcal{P}^{(2)}_{i,\ell=1}+\mathcal{P}^{(2)}_{i,\ell=2} \, .
\end{equation}
%The general expression for the change in the eigenvalue of the atomic Hamiltonian between the channels with different orbital angular momentum is
%\begin{equation}     \sum_{ij} \big[  { H } _{\text{at}} \, ,\, \mathcal{P}^{(n)}_{i,\ell} \, h_i^\dagger  \tilde{\mathbf{T}}_{ij} h_j^{\phantom{\dagger}} \, \mathcal{P}^{(n')}_{j,\ell'} \big] = \Delta V (n \ell ; n' \ell' ) \sum_{ij}  \mathcal{P}^{(n)}_{i,\ell} \, h_i^\dagger  \tilde{\mathbf{T}}_{ij} h_j^{\phantom{\dagger}} \, \mathcal{P}^{(n')}_{j,\ell'} \;. \end{equation}
The relevant commutator to cancel the first-order terms is
\begin{equation}
    \sum_{ij} \big[  { H } _{\text{at}} \, ,\, \mathcal{P}^{(2)}_{i,\ell} \, h_i^\dagger  \tilde{\mathbf{T}}_{ij} h_j^{\phantom{\dagger}} \, \mathcal{P}^{(1)}_{j} \big] = \Delta E_{\ell}  \, \sum_{ij}  \mathcal{P}^{(2)}_{i,\ell} \, h_i^\dagger  \tilde{\mathbf{T}}_{ij} h_j^{\phantom{\dagger}} \, \mathcal{P}^{(1)}_{j} \, , \label{eq:II-1-comm}
\end{equation}
where
\begin{equation}
    \Delta E_0 = U + 2 J_H \, , \quad
    \Delta E_1 = U - 3 J_H \, , \quad
    \Delta E_2 = U - J_H \, . \label{eq:II-1-delta-E}
\end{equation}



%We can use a canonical transformation,
%\begin{equation}    H \to \tilde{ { H } } = e^S H e^{-S} = H + [S,H]+\frac{1}{2}[S,[S,H]] + \cdots,\end{equation}
%which eliminates terms that change $\bar{N}$ up to order $(t/U)^2$.
Lets write the canonical transformation as $S=S_1+S_2+...$ where $S_k$ is of order $k$ in the perturbation, i.e. $S_k\propto (t/U)^k$.
The transformed  Hamiltonian \eqref{eq:II-1-S-transf} is
\begin{equation}
\begin{split}
        \tilde{ { H } } &=  T_0 +  { H } _{\text{at}} + \left( T_1 + T_{-1} + [S_1,  { H } _{\text{at}}] \right) +  \\[4pt]
    &+ [S_1, T_0+T_1 + T_{-1}] + \frac{1}{2} [S_1,[S_1,  { H } _{\text{at}}]] + [S_2,  { H } _{\text{at}}] + \cdots \; ,
\end{split}
\end{equation}
where the terms of order $(t/U)^2$ are in the second line.


The goal of the canonical transformation is to choose $S$ such that the perturbation is canceled in the first order.
To eliminate the perturbation\footnote{Note that we do not include the term $T_0$ of $ { H } _{\text{hop}}$ in the `perturbation', even though it is a term $\propto t$, as it does not change the occupation.  } in the first order, we must choose $[S_1,  { H } _{\text{at}}] = - T_1 - T_{-1}$ which is equivalent to
\begin{align}
    S_1 &= S_1^{+} - S_1^{-} \; ,  \quad
    (S_1^{-})^\dagger = S_1^{+} \; , \quad
    [ { H } _{\text{at}},S_1^{+}] =   T_{ 1 } \; .
\end{align}
Then, choose
\begin{equation}
    S_1^{+}  = \sum_{ij}\sum_{\ell} \frac{1}{\Delta E_{\ell}} \mathcal{P}^{(2)}_{i,\ell} \, h_i^\dagger  \tilde{\mathbf{T}}_{ij} h_j^{\phantom{\dagger}} \, \mathcal{P}^{(1)}_{j}  \; .
\end{equation}
We find the Hamiltonian,
\begin{equation}
    \begin{split}
        \tilde{ { H } } &=   T_0 +  { H } _{\text{at}} +  [S_1, T_0]  + \frac{1}{2} [S_1, T_1 + T_{-1}] + [S_2,  { H } _{\text{at}}] + \cdots \; .
\end{split}
\end{equation}
To simplify the Hamiltonian we also choose
\begin{equation}
     [S_1, T_0]  +
     [S_2,  { H } _{\text{at}}] = 0 \; . \label{eq:II-1-S2}
\end{equation}
To compute the effective Hamiltonian $\tilde{H}$ in second-order perturbation theory one does not need to compute $S_2$ as it gets canceled by the relation \eqref{eq:II-1-S2}. However, for computing a transformed observable
$\tilde{\mathcal{O}} = e^{\, S}  \;  \mathcal{O} \, e^{-S}$ in the same or higher order of perturbation $S_2$ is needed.%  in second order of the perturbation.

\subsection{Effective Hamiltonian}
The second step in the process to get the effective Hamiltonian is to restrict to the low-energy subspace of the unperturbed Hamiltonian.
Let us define the projection operator to the low-energy sector of one hole per site
\begin{equation}
    \mathcal{P}_{1\mathrm{h}} =
    \prod_i \mathcal{P}_{i}^{(1)} \; , \quad  \mathcal{P}_{i}^{(1)} = \sum_{s_i} \ket{s_i}\bra{s_i}\; .
\end{equation}
In this subspace we have
\begin{equation}
    \mathcal{P}_{1\mathrm{h}} V_N \mathcal{P}_{1\mathrm{h}} = 0 \, , \quad
    \mathcal{P}_{1\mathrm{h}} V_S \mathcal{P}_{1\mathrm{h}} = - 2J_H \, \frac{3}{4} \bar{N}   \, , \quad
    \mathcal{P}_{1\mathrm{h}} V_L \mathcal{P}_{1\mathrm{h}} = -  \frac{J_H}{2} \, 2 \, \bar{N} \, ,
\end{equation}
where $\bar{N} = \sum_i 1$ is the total number of holes. Then we conclude that the projection of the atomic Hamiltonian is a constant %\textcolor{red!70}{[factor of 2?]}
$\mathcal{P}_{1\mathrm{h}}  { H } _{\text{at}} \mathcal{P}_{1\mathrm{h}} = - 5J_H\bar{N} /2$.
The projection of the hopping term in this subspace is zero $\mathcal{P}_{1\mathrm{h}} \,  { H } _{\text{hop}}\,\mathcal{P}_{1\mathrm{h}} = 0$ and the commutator $[S_1, T_0]$ also has zero projection since the terms $S_1^{\pm}T_0$ and $T_0S_1^{\pm}$ go out of the subspace with one hole per site. For the commutator $[S_1, T_{\pm 1}]$, the terms that stay in this subspace are  $[S_1^+, T_{- 1}]$ and  $[S_1^-, T_{+1}]$.


Projecting the Hamiltonian in the low-energy subspace and subtracting a constant term $\sim J_H\bar{N}$, and ignoring the term with $S_2$ we then have
\begin{align}\hspace*{-15mm}
    \mathcal{P}_{1\mathrm{h}} \tilde{ { H } }  \mathcal{P}_{1\mathrm{h}}
&=
     \mathcal{P}_{1\mathrm{h}}  \Big(\, \frac{1}{2} [S_1, T_1 + T_{-1}] \, \Big)
     \mathcal{P}_{1\mathrm{h}} \, , \\[6pt]
     &=
    \frac{1}{2} \mathcal{P}_{1\mathrm{h}}  \Big( \,  [S_1^+, T_{- 1}] - [S_1^-, T_{+1}] \, \Big) \mathcal{P}_{1\mathrm{h}} \, , \\[6pt]
&=  \frac{1}{2}
     \mathcal{P}_{1\mathrm{h}}  [S_1^{(+)}, T_{-1}]  \mathcal{P}_{1\mathrm{h}} + \text{h.c.}\\[6pt]
     &=- \frac{1}{2}
     \mathcal{P}_{1\mathrm{h}} T_{-1} S_1^{(+)}  \mathcal{P}_{1\mathrm{h}} + \text{h.c.} \, .
\end{align}
In the last line, note that $S^+$ maps into the sector with one double occupancy, thus  $T_{-1} S_1^{(+)}$  has non-zero projection while $\mathcal{P}_{1\mathrm{h}}  S_1^{(+)} T_{-1}  \mathcal{P}_{1\mathrm{h}} =0$.
%[In Rodrigo's notes, he said that the we must consider only the term with double occupancy $[S_1^{(+)}, T_{-1}] \to  - T_{-1} S_1^{(+)} $. This is the term that will produce the spin Hamiltonian. What is the term that we ignore?]




The effective Hamiltonian in the low-energy subspace is
\begin{equation}\hspace*{-10mm}
    \begin{split}
        \mathcal{P}_{1\mathrm{h}} \tilde{ { H } }  \mathcal{P}_{1\mathrm{h}}
      &= -
      \sum_{ij}\sum_{\ell} \frac{1}{\Delta E_{\ell}} \mathcal{P}_{1\mathrm{h}}  \, \big( h_j^\dagger  \tilde{\mathbf{T}}_{ji} h_i^{\phantom{\dagger}} \big) \, \mathcal{P}^{(2)}_{i,\ell} \, \big( h_i^\dagger  \tilde{\mathbf{T}}_{ij} h_j^{\phantom{\dagger}} \big)\, \mathcal{P}_{1\mathrm{h}} \; .
    \end{split}
\end{equation}
Considering the matrix elements between a pair of sites $i,j$. Let us write $\tilde{ { H } }=\sum_{\ell} \tilde{ { H } }_{\ell}$, and consider the contribution from each channel separately
\end{comment}

\section{Hamiltonian in spin-orbital basis}

The matrix elements of the Hamiltonian \eqref{micro-Heff} in the basis of spin-orbital basis $\ket{s}=\ket{a,\sigma}$ are given by
\begin{adjustwidth}{0mm}{-5mm}
\begin{equation}\hspace*{-5mm}
    \begin{split}
        \hspace*{-5mm}\bra{s_i', s_j'} \tilde{ { H } }_{\ell}  \ket{s_i, s_j}
      &= -   \frac{1}{\Delta E_{\ell}} \bra{s_i', s_j'}  \, \big( h_j^\dagger  \tilde{\mathbf{T}}_{ji} h_i^{\phantom{\dagger}} \big) \, \mathcal{P}^{(2)}_{i,\ell} \, \big( h_i^\dagger  \tilde{\mathbf{T}}_{ij} h_j^{\phantom{\dagger}} \big)\, \ket{s_i, s_j} \;
      \\
      &= -   \frac{1}{\Delta E_{\ell}} \sum_{s_1,s_2,s_3,s_4}\big(\tilde{\mathbf{T}}_{ji} \big)_{s_1,s_2}\big(  \tilde{\mathbf{T}}_{ij}\big)_{s_3,s_4}   \bra{s_i', s_j'}  \,  h_{js_1}^\dagger   h_{is_2}^{\phantom{\dagger}} \, \mathcal{P}^{(2)}_{i,\ell} \, h_{is_3}^\dagger h_{js_4}^{\phantom{\dagger}} \, \ket{s_i, s_j}  \\
      &= -   \frac{1}{\Delta E_{\ell}} \sum_{s_1,s_2,s_3,s_4}\big(\tilde{\mathbf{T}}_{ji} \big)_{s_1,s_2}\big(  \tilde{\mathbf{T}}_{ij}\big)_{s_3,s_4}
      \bra{s_i'}    h_{is_2}^{\phantom{\dagger}} \mathcal{P}^{(2)}_{i,\ell}  h_{is_3}^\dagger  \ket{s_i}
      \bra{s_j'}    h_{js_1}^\dagger   h_{js_4}^{\phantom{\dagger}} \ket{s_j}
   ,
    \end{split}
\end{equation}\end{adjustwidth}
where we separated the contribution for each sector $\ell$.
The matrix element at site $j$ are simply, $
    \bra{s_j'}    h_{js_1}^\dagger   h_{js_4}^{\phantom{\dagger}} \ket{s_j}  =
    \delta_{s_4,s_j}\delta_{s_1,s_j' } $, and the expression reduces to
\begin{equation}
    \begin{split}\hspace*{-8mm}
        \bra{s_i', s_j'}\! \tilde{ { H } }_{\ell} \!\ket{s_i, s_j} \!
      &= -   \frac{1}{\Delta E_{\ell}} \!
      \sum_{s_2,s_3}\!\!\big({\mathbf{T}}_{ji} \big)_{ {a}_j', {a}_2}\!\big(  {\mathbf{T}}_{ij}\big)_{ {a}_3,{a}_j}\!\delta_{\sigma_j'\sigma_2} \delta_{\sigma_3\sigma_j}
      \bra{s_i'}    h_{is_2}^{\phantom{\dagger}} \mathcal{P}^{(2)}_{i,\ell}  h_{is_3}^\dagger  \ket{s_i}                .\hspace*{-1mm}
    \end{split}
\end{equation}
Let us consider each channel separately to calculate the matrix elements at site $i$.
\begin{align}
   %\begin{split}      \mathcal{P}^{(2)}_{i,\ell=0}  &= \ket{\Psi_{i,\ell=0}^{(2)}}\bra{\Psi_{i,\ell=0}^{(2)}} \; ,  \\[6pt]
    %\ket{\Psi_{i,\ell=0}^{(2)}} &= \frac{1}{\sqrt{3}} \sum_a h_{ia\uparrow}^\dagger h_{ia\downarrow}^{\dagger \phantom{\dagger}}\ket{F} \; .   \end{split}\\
   \bra{s_i'}    h_{is_2}^{\phantom{\dagger}} \mathcal{P}^{(2)}_{i,\ell=0}  h_{is_3}^\dagger  \ket{s_i}
    &=
    \frac{1}{3} \sigma_2\sigma_3 \delta_{a_{i}'a_{2}}
    \Big( \delta_{\bar\sigma_i', \sigma_2} \Big)
    \delta_{a_{i}a_{3}}
    \Big( \delta_{\bar\sigma_i,\sigma_3} \Big)
        \\[6pt]
\bra{s_i'}    h_{is_2}^{\phantom{\dagger}} \mathcal{P}^{(2)}_{i,\ell=1}  h_{is_3}^\dagger  \ket{s_i}
&=
\frac{1}{2}
\Big( \delta_{a_{i}'a_{i}} \delta_{a_{2}a_{3}} - (\alpha_3 \leftrightarrow \alpha_i)  \Big)
\Big( \delta_{\sigma_{i}'\sigma_{i}} \delta_{\sigma_{2}\sigma_{3}} + (\sigma_3 \leftrightarrow \sigma_i)  \Big)
\\[6pt]
    \begin{split}
      \bra{s_i'}    h_{is_2}^{\phantom{\dagger}} \mathcal{P}^{(2)}_{i,\ell=2}  h_{is_3}^\dagger  \ket{s_i}
     &=
     \bra{s_i'}    h_{is_2}^{\phantom{\dagger}}    h_{is_3}^\dagger  \ket{s_i} \\
     &
     -
     \bra{s_i'}    h_{is_2}^{\phantom{\dagger}} \mathcal{P}^{(2)}_{i,\ell=0}  h_{is_3}^\dagger  \ket{s_i}
     -
     \bra{s_i'}    h_{is_2}^{\phantom{\dagger}} \mathcal{P}^{(2)}_{i,\ell=1}  h_{is_3}^\dagger  \ket{s_i}
     \; .
    \end{split}
\end{align}
The final expression for the matrix elements of the Hamiltonian is
\begin{align}
     \bra{s_i', s_j'} \tilde{ { H } }_0  \ket{s_i, s_j}
    &= -   \frac{1}{3 \Delta E_0}
    \big({\mathbf{T}}_{ji} \big)_{a_j'\,a_i'} \big(  {\mathbf{T}}_{ij}\big)_{a_i\,a_j} \sigma_j'\sigma_j      \delta_{\bar\sigma_i'\, \sigma_j'}
    \delta_{\bar\sigma_i\,\sigma_j}
    \; , \\
    \begin{split}
    \bra{s_i', s_j'} \tilde{ { H } }_1  \ket{s_i, s_j}
    &= -   \frac{1}{2 \Delta E_1} \Big[
    \big({\mathbf{T}}_{ji}   {\mathbf{T}}_{ij}\big)_{a_j'\,a_j}  \delta_{a_i'\,a_i}
    - \big({\mathbf{T}}_{ji} \big)_{a_j'\,a_i} \big(  {\mathbf{T}}_{ij}\big)_{a_i'\,a_j} \Big] \;\\ &\times  \big( \,
      \delta_{\sigma_j'\, \sigma_j}  \delta_{\sigma_i\,\sigma_i'}  +
      \delta_{\sigma_j'\, \sigma_i}  \delta_{\sigma_j\,\sigma_i'} \,\big)  \; .
    \end{split}\\
    \begin{split}
        \bra{s_i', s_j'} \tilde{ { H } }_2  \ket{s_i, s_j}
      &= -
      \frac{1}{\Delta E_2} \!      \Big[
       \big({\mathbf{T}}_{ji}  {\mathbf{T}}_{ij}\big)_{a_j'\,a_j}\!\delta_{\sigma_j'\sigma_j} \delta_{s_i' s_i}       -
      \big({\mathbf{T}}_{ji} \big)_{a_j'\,a_i}\big(  {\mathbf{T}}_{ij}\big)_{a_i'\,a_j}\!\delta_{\sigma_j'\,\sigma_i} \delta_{\sigma_i'\,\sigma_j}        \Big]      \\
      & -
      \frac{\Delta E_0}{\Delta E_2} \bra{s_i', s_j'} \tilde{ { H } }_0  \ket{s_i, s_j}
      -
      \frac{\Delta E_1}{\Delta E_2} \bra{s_i', s_j'} \tilde{ { H } }_1  \ket{s_i, s_j}            .
    \end{split}
\end{align}
These matrix elements are expressed in the full spin-orbital basis, in the next section we will compute the representation in the effective $j=1/2$ basis.


\section{Effective spin-half Hamiltonian}


The matrix in the $j_{\text{eff}}=1/2$ subspace is obtained by the projection $\mathcal{P}_{\frac{1}{2}}$.
To go from the basis $\{ \ket{s} \} =\big\{ \ket{x\uparrow},  \ket{x\downarrow}, \; \,\ket{y\uparrow},  \ket{y\downarrow}, \; \,\ket{z\uparrow},  \ket{z\downarrow} \big\}$ to the basis spin-orbital-half is a change of basis and projection in the subspace spanned by the vectors $\{\ket{+},\ket{-}\}$. Algebraically, we need to compute matrix multiplication in eq.~\eqref{eq:oij}. This process is implemented symbolically using \texttt{Mathematica}.

As an example of this projection, consider a magnetic field $B_z$ applied in the $z$ direction. The Hamiltonian for this operator is $ { H } _{h}=-B_z S^z=\frac{1}{2}\sum B_z \, h_i^\dagger \big(\mathbb{1}_3\otimes\sigma_z\big)h_i$. The matrix for this term is $\bra{s_i'}    { H } _{Z} \ket{s_i} =  \frac{1}{2}B_z \,   \big(\mathbb{1}_3\otimes\sigma_z\big)_{s_i',s_i}$, which after the projection gives
$\mathcal{P}_{\frac{1}{2}}  { H } _{Z}  \mathcal{P}_{\frac{1}{2}} =  \frac{1}{2}B_z \,   \mathbb{M}^\dagger \big(\mathbb{1}_3\otimes\sigma_z\big) \mathbb{M} = \frac{1}{6} B_z \big( \sigma_x\sigma_z\sigma_x+\sigma_y\sigma_z\sigma_y+\sigma_z\sigma_z\sigma_z \big) = - \frac{1}{6} B_z \sigma_z  $.

The Hamiltonian $\tilde{ { H } }$ is a sum of terms that act on two sites.
After the projection into the effective spin-1/2 space we find (up to constants)
\begin{align}
    \mathcal{P}_{\frac{1}{2}}\tilde{ { H } }_0   \mathcal{P}_{\frac{1}{2}}
    &=
    \frac{(2  t_{1}+ t_{3})^2}{27 \, \Delta E_0}
    \sum_{\left<ij\right>_\gamma } \, \bm{\sigma}_i\cdot\bm{\sigma}_j
    \;,\\
    \begin{split}
        \mathcal{P}_{\frac{1}{2}}\tilde{ { H } }_1   \mathcal{P}_{\frac{1}{2}}
    &=
    \frac{1}{9\Delta E_1} \sum_{\gamma(\alpha\beta)}
    \sum_{\left<ij\right>_\gamma } \Bigg[
    2t_{1} ( t_{1}+2  t_{3}) \bm{\sigma}_i\cdot\bm{\sigma}_j
    \\ &+  \Big(\big( t_{1}- t_{3}\big)^2-3t_{2}^2 \Big) \sigma_i^\gamma\sigma_j^\gamma
    + 2t_{2} ( t_{1}- t_{3})    \big(\sigma_i^\alpha \sigma_j^\beta +\sigma_i^\beta \sigma_j^\alpha  \big)  \Bigg]
    \; ,
    \end{split} \\
    \mathcal{P}_{\frac{1}{2}}\tilde{ { H } }_2   \mathcal{P}_{\frac{1}{2}}
    &=
    \frac{(2  t_{1}+ t_{3})^2}{9 \, \Delta E_2}
    \sum_{\left<ij\right>_\gamma }  \bm{\sigma}_i\cdot\bm{\sigma}_j
    -\frac{\Delta E_0}{\Delta E_2} \mathcal{P}_{\frac{1}{2}}\tilde{ { H } }_0   \mathcal{P}_{\frac{1}{2}}
    -\frac{\Delta E_1}{\Delta E_2}\mathcal{P}_{\frac{1}{2}}\tilde{ { H } }_1   \mathcal{P}_{\frac{1}{2}}
\end{align}
Summing all these terms and grouping like terms,
one finds the extended Kitaev Hamiltonian in eq. \eqref{eq:ext-Ham} of the main text.


\section{Second order canonical transformation}
The second order $S_2$ in the canonical transformation expansion is determined by the condition in eq.\eqref{eq:II-1-S2}, which I will write again for convenience
\begin{equation}
    \begin{split}
       0 &= [S_1, T_0]  +
     [S_2, H_{\text{at}}]
     \; , \\[6pt]
     S_1  &= \sum_{ij}\sum_{\ell} \frac{1}{\Delta E_{\ell}} \mathcal{P}^{(2)}_{i,\ell} \, h_i^\dagger  \tilde{\mathbf{T}}_{ij} h_j^{\phantom{\dagger}} \, \mathcal{P}^{(1)}_{j} \;  - \mathrm{h.c}
     \; , \\
     H_{\mathrm{at}} &=\sum_i \left( \frac{U-3J_H}{2}(N_i-5)^2 - 2J_H \mathbf{S}_i^2 - \frac{J_H}{2} \mathbf{L}_i^2 \right)
     \; ,\\
     T_0 &=  - \sum_{ij} \Big( \cP_i^{(1)} \;  h_i^\dagger  \tilde{\mathbf{T}}_{ij} h_j^{\phantom{\dagger}}  \, \cP_j^{(1)}
     +
     \cP_i^{(2)} \;  h_i^\dagger  \tilde{\mathbf{T}}_{ij} h_j^{\phantom{\dagger}}  \, \cP_j^{(2)} \Big) \; .
    \end{split}
\end{equation}
To determine the effective Hamiltonian in the second order, we do not need to know $S_2$. It is enough to know that $S_2$ satisfies the previous equation. However, we need $S_2$ to determine the charge density in the third order.

Using $\Delta V=\Delta V_{N}+\Delta V_{S}+\Delta V_L$ to define the change in the eigenvalue of the atomic Hamiltonian $H_{\mathrm{at}}$ by hopping off one hole. Consider $\Delta V = \Delta V(n,\ell; n',\ell')$ corresponding to hopping with a configuration with $n'$ holes with angular momentum $\ell'$ to a configuration with $n$ holes and angular momentum $\ell$. This expression is defined by
\begin{equation}
    \sum_{ij} \; \left[ H_{\mathrm{at}} \; , \; \cP_{i,\ell}^{(n)} \, h_i^\dagger   \tilde{\mb{T}}_{ij} h_j^\pd \, \cP_{j,\ell'}^{(n')} \right]
    =
    \sum_{ij} \;    \Delta V(n,\ell; n',\ell')
    \; \cP_{i,\ell}^{(n)} \, h_i^\dagger  \mb{T}_{ij}  h_j^\pd \, \cP_{j,\ell'}^{(n')}
\end{equation}
Where we are interested in $n,n' \in \{ 0,1,2 \}$. For one hole only, only $\ell=1$ is allowed, and these energies change by adding one hole to a configuration where there was already one hole are $\Delta E_\ell$ in eq.\eqref{eq:II-1-delta-E}.
The values for these changes in the potential are
\begin{align}
    \Delta V(1,\ell; 1,\ell') &= 0 \; , \\
    \Delta V(2,\ell; 1,\ell') &= \delta_{\ell',1}  \; \Delta E_{\ell} \; , \\
    \Delta V(1,\ell; 2,\ell') &= -\delta_{\ell,1}  \; \Delta E_{\ell'} \; , \\
    \Delta V(2,\ell; 2,\ell') &= \Delta E_{\ell} - \Delta R_{\ell'} \; .
\end{align}
The expression with $n=0$ or $n'=0$ does not appear in intermediate steps when projecting the initial and final state into the one hole per site subspace.

The expression for the canonical transformation is
\begin{equation}
    S_2 =    \sum_{ijkl}\!\sum_{\ell\,\ell'\ell''}\!\sum_{n=1,2}\!
    A(\ell;n,\ell',\ell'')
    \left[
    \mathcal{P}^{(2)}_{i,\ell} \, h_i^\dagger  \tilde{\mathbf{T}}_{ij} h_j^{\phantom{\dagger}}  \mathcal{P}^{(1)}_{j}  ,\cP_{k,\ell'}^{(n)} \; h_k^\dagger  \tilde{\mb{T}}_{kl}  h_l^\pd  \cP_{l,\ell''}^{(n)}
    \right]
    - \mathrm{h.c} \; ,  \label{eq:S2}
\end{equation}
with coefficient
\begin{equation}
    A(\ell;n,\ell',\ell'') =  \frac{1}{\Delta E_\ell \; \left( \Delta E_{\ell'} + \Delta V(n,\ell';n,\ell'') \; \right)} \; .
\end{equation}
Some of the processes contributing to $S_2$ are: one hole at $j$ hops to $i$ forming a state with angular momentum $\ell$ then one of these doubly occupied holes hops to another site. Many other processes are happening in $S_2$, however, the projection into the subspace with one hole per site will simplify the expressions.







\section{Charge density matrix elements}
After projecting in the $j_{\text{eff}}=1/2$ subspace, the general expression for the charge density is \eqref{spincorr} with coefficients given by
\begin{align}
    \mathcal{C}_{jkl}^{\mu\nu\rho} &=
    \frac{1}{8} \; \text{Tr} \left( \cP_{\frac12}  \delta n_{l,(jk)}^{(3)} \cP_{\frac12} \sigma^\mu_j \sigma^\nu_k \sigma^\rho_l  \right) \; , \\[6pt]
    \begin{split}
     &=
    \frac{1}{8} \sum_{\sigma_j',\sigma_k',\sigma_l'}\sum_{\sigma_j, \sigma_k, \sigma_l} \\
    & \;\,
    \braket{\sigma_j',\sigma_k',\sigma_l' | \cP_{\frac12}  \delta n_l \cP_{\frac12}  | \sigma_j, \sigma_k, \sigma_l}
    \big(\sigma^\mu\big)_{\sigma_j,\sigma_j'}
    \big(\sigma^\nu\big)_{\sigma_k,\sigma_k'}
    \big(\sigma^\rho\big)_{\sigma_l,\sigma_l'}
     \; , \!
    \end{split}
\end{align}
where $\big(\sigma^\mu\big)_{\sigma,\sigma'}$ are the matrix elements of Pauli matrices $(\sigma^\mu)=(\mathbb{1}_2,\sigma^x,\sigma^y,\sigma^z)$ and the sum runs over the effective spin-half indices $\sigma=+,-$.

The matrix elements of the charge density for points in a triangle $ljk$ are
\begin{equation}
\begin{split}
    \braket{\sigma_j',\sigma_k',\sigma_l' | \cP_{\frac12}
 \delta n_l \cP_{\frac12}  | \sigma_j, \sigma_k, \sigma_l}
 &=\!
 \sum_{s_j's_k's_l',s_js_ks_l}\!
 \mathbb{M}_{\sigma_j' \, s_j'}^\dagger   \mathbb{M}_{\sigma_k' \, s_k'}^\dagger
 \mathbb{M}_{\sigma_l' \, s_l'}^\dagger
 \mathbb{M}_{s_j\, \sigma_j}
 \mathbb{M}_{s_k\, \sigma_k}
 \mathbb{M}_{s_l\, \sigma_l}  \\
 &\phantom{\sum_{s_j's_k's_l',s_js_ks_l}}
 \braket{s_j',s_k',s_l' | \delta n_l | s_j, s_k, s_l}
\end{split}
\end{equation}
where the matrix elements in the spin-orbital space are
\begin{align}
    \hspace*{-9mm}
        \braket{s_j',s_k',s_l' | \delta n_l | s_j, s_k, s_l} &=\!
        \braket{s_j',s_k',s_l' | \delta n_{l,jk} | s_j, s_k, s_l}+
        \braket{s_j',s_k',s_l' | \delta n_{l,kj} | s_j, s_k, s_l} \; , \\
        &=:\!
        \braket{s_j',s_k',s_l' | \delta n_{l,(jk)} | s_j, s_k, s_l}
\end{align}
for which
\begin{equation}
    \begin{split}\hspace*{-9mm}
        &\braket{s_j',s_k',s_l' | \delta n_{l,jk} | s_j, s_k, s_l} =\\
        \hspace*{-9mm}        &=\!
        -\!\!\!\!\sum_{s_2,s_3,s_4,s_5} \!\!\!\!
    \big({\mathbf{T}}_{jk} \big)_{ {a}_j', {a}_2} \big({\mathbf{T}}_{kl} \big)_{ {a}_3, {a}_4}\big({\mathbf{T}}_{lj} \big)_{ {a}_5, {a}_j}
    \sum_{\ell, \ell'}\frac{F_{\ell'}(s_k',s_2,s_3,s_k) \; F_{\ell}(s_l',s_4,s_5,s_l) }{ (\Delta E_\ell)^2 \; \Delta E_{ \ell'} }
    \\
    &+\!\!\!\!\sum_{s_2,s_3,s_4,s_5} \!\!\!\!
    \big({\mathbf{T}}_{lk} \big)_{ {a}_j', {a}_2} \big({\mathbf{T}}_{kj} \big)_{ {a}_3, {a}_4}\big({\mathbf{T}}_{jl} \big)_{ {a}_5, {a}_l}
    \sum_{\ell, \ell'}\frac{F_{\ell}(s_j',s_4,s_5,s_j) \; F_{\ell'}(s_k',s_2,s_3,s_k) }{ (\Delta E_\ell)^2 \; \Delta E_{ \ell'} } \\
    & -\!\sum_{s_2,s_5} \!
    \big({\mathbf{T}}_{jk} \big)_{ {a}_k', {a}_2} \big({\mathbf{T}}_{kl} \big)_{ {a}_l', {a}_k}\big({\mathbf{T}}_{lj} \big)_{ {a}_5, {a}_l}
    \sum_{\ell}\frac{F_{\ell}(s_j',s_2,s_5,s_j)  }{ (\Delta E_\ell)^3 }
    \\
    &+\!\sum_{s_2,s_5} \!
    \big({\mathbf{T}}_{lk} \big)_{ {a}_k', {a}_2} \big({\mathbf{T}}_{kj} \big)_{ {a}_j', {a}_k}\big({\mathbf{T}}_{jl} \big)_{ {a}_5, {a}_j}
    \sum_{\ell}\frac{F_{\ell}(s_l',s_2,s_5,s_l) \;  }{ (\Delta E_\ell)^3 } + \text{h.c.}
    \; ,  \label{eq:II-5-dcharge-jk}
    \end{split}
\end{equation}
where
\begin{equation}
    F_{\ell}(s_1,s_2,s_3,s_4) = \braket{ s_1 | h_{s_2}^{\phantom{\dagger}} \; \mathcal{P}_{\ell}^{(2)} \; h_{s_3}^{\dagger} | s_4 }
\end{equation}
and $\Delta E_\ell$ are defined in eq. \eqref{eq:II-1-delta-E}.
Note that the matrices $\mathbf{T}_{ij}$ are only non-zero between sites $ij$ that are either nearest-neighbors or next-nearest-neighbors.
Importantly, when exchanging $j$ and $k$, the matrices $\mathbf{T}_{jk}, \mathbf{T}_{kl}, \mathbf{T}_{lj}$ change, as the matrixes are different for each type of bonds\footnote{Bonds in a generalized sense to include next-nearest-neighbors bonds.} between $(jk)$, $(kl)$ and $(lj)$.


The coefficients $\mathcal{C}_{jkl}$ are invariant by translation of the honeycomb lattice, so it only depends on the relative position between the sites $jkl$. These positions can be any three points in the lattice, but the coefficients are independent of the position of $l$ and only differ by the relative position of $j$ and $k$.
Thus, let us use the relative orientation of these points to label the set of coefficients. There are $s=1, ..., 18$ ways to arrange three points in the honeycomb lattice such that there are two pairs of NN bonds and the other pair is in a NNN bond, these are displayed in figure \ref{fig:AB_triangles}.
Where $s=1, ..., 9$ refer to when $l$ is in $A$ sublattice, and $s=10, ..., 18$ for $l$ in the $B$ sublattice. The position dependence in the coefficients $\mathcal{C}_{(s)}^{\mu\nu\rho}=\mathcal{C}_{jkl}^{\mu\nu\rho}$ are fully described by the type of triangle $s$.

The charge density at site $l$  is the sum of the contribution of $9$ triangles%, and for each triangle indexed by $s$ the coefficients only depend on the type of the triangle
\begin{adjustwidth}{-0mm}{-0mm}\vspace*{-5mm}
\begin{equation}
\begin{split}
        \braket{\delta n_l} &= \sum_s \delta n_{l,(s)} \; ,
        \\
    \delta n_{l,(s)} &= \sum_{jk}
    \sum_{\alpha\beta} \Big(  \mathcal{C}_{(s)}^{\alpha 0 \beta } \braket{\sigma_j^\alpha \sigma_l^\beta} +\mathcal{C}_{(s)}^{ 0 \alpha\beta } \braket{\sigma_k^\alpha \sigma_l^\beta} +
         \mathcal{C}_{(s)}^{   \alpha\beta 0 } \braket{\sigma_j^\alpha \sigma_k^\beta}  \Big) \; , \label{eq:II-5-density}
\end{split}
\end{equation}\end{adjustwidth}
where the sum in $jk$ is such that $(jkl)$ is \textit{the} triangle of the form $(s)$.
Now we need to compute the coefficients $ \mathcal{C}_{(s)}^{\mu\nu\rho} $ for each triangle $s$.
%\paragraph{First type of configuration}
For example, let us consider the triangle in figure \ref{fig:triang-jkl}.
The label for this configuration is $s=11$.
This triangle is defined by the sites $(jlk)$ where $jl$ is in an $x$ bond, $lk$ is in a $y$-bond and the NNN bond between $jk$ is of type $z$. Thus,
the hopping matrices for this process are
\begin{equation}
    \mathbf{T}_{jl}  =    \mathbf{T}_{\braket{x}} ,    \;
    \mathbf{T}_{lk}  =    \mathbf{T}_{\braket{y}} , \;
    \mathbf{T}_{jk}  =    \mathbf{T}_{\braket{\braket{z}}}  .
\end{equation}
These are the hopping matrices for the configuration with $j$ and $k$ as in figure \ref{fig:triang-jkl}. However, in the charge density, $\delta n_{l (jk)}$ we are also considering the contribution due to swapped positions $j \leftrightarrow k$. The hopping matrices for this other contribution are
\begin{equation}
    \mathbf{T}_{jl}  =    \mathbf{T}_{\braket{y}} ,    \;
    \mathbf{T}_{lk}  =    \mathbf{T}_{\braket{x}} , \;
    \mathbf{T}_{jk}  =    \mathbf{T}_{\braket{\braket{z}}}  .
\end{equation}
Using these matrices in equation \eqref{eq:II-5-dcharge-jk} we can compute $\braket{s_j',s_k',s_l' | \delta n_{l,jk} | s_j, s_k, s_l}$ and
$\braket{s_j',s_k',s_l' | \delta n_{l,kj} | s_j, s_k, s_l}$. Then, using equation \eqref{eq:II-5-C} we can compute the coefficients $\mathcal{C}_{(s)}^{\mu\nu\rho}=\mathcal{C}_{jkl}^{\mu\nu\rho}$ for this configuration $s=11$.

In principle we should compute the coefficients $\mathcal{C}_{(s)}^{\mu\nu\rho}$ for each set of triangles $(s)$.
However one can see that they are not all independent.
For example in fig. \ref{fig:AB_triangles} we see that the triangle $s=12$ is related by a $120^\circ$ rotation from the $s=11$ triangle.
Identifying the bonds $xyz \overset{\pi}{\longrightarrow} yzx$  one gets the relation between the coefficients
$\mathcal{C}_{(s=12)}^{\pi(\mu)\,\pi(\nu)\,\pi(\rho)}=\mathcal{C}_{(s=11)}^{\mu\nu\rho}$.

In total, we need to compute eq. \eqref{eq:II-5-dcharge-jk} and then project each of the $64$ entries (of which at most $27$ are non-zero) of $C_{(s)}$ for  each of the $18$ triangles. This type of repetitive task claims for an automated process, which was implemented using symbolic computation in \texttt{Mathematica}. %The computation of the $27$
Here I will not display all these coefficients, but I can write for instance the dominant coefficient close to the Kitaev limit  ($\eta=J_H/U$)
\begin{equation}\hspace*{-4mm}
    \max_{ }  \big| \mathcal{C}^{\mu\nu\rho}_{(s)}  \big|  =
    \big| \mathcal{C}^{x0x}_{(11)}  \big|  = \frac{t_2^2  t_2^\prime}{U^3} \frac{\eta \left(6-23\eta+22\eta^2\right)}{\; 27 \left(1-\eta\right)^3\left(1-3\eta\right)^3 \;} \approx 7 \times 10^{-6} \; .
\end{equation}
%Beyond the $3$ triangles related by $120^\circ$ rotation of the triangle in fig  \ref{fig:triang-jkl} (b), there are more $6$ triangles. These other $6$ are also related by some symmetry.Now consider another configuration that contributes to the charge density at $l$.  Let $j$ and $k$ be such that $\braket{jl}_x$ and $\braket{jk}_y$.
%The charge density is  order
The overall proportional coefficient is the expected $J_H\,t^2  t^\prime/U^4$. This coefficient correctly reproduces the coefficient for the Kitaev limit on Ref. \cite{Pereira2020}.

\begin{figure}[t]
    \centering
    \includegraphics[width=0.49\textwidth]{Chapters/Appendix/Figures/micro/A_triangle.pdf}\hspace{1mm}%
    \includegraphics[width=0.49\textwidth]{Chapters/Appendix/Figures/micro/B_triangle.pdf}
    \caption{Triangles configurations. For $s=$01-09 are the triangles where $l$ is in the $A$ sublattice, whereas for $s=$10-18 $l$ is in the $B$ sublattice. The triangles are label by the sublattice $\sigma=A,B$ of $l$, and the relative position $(m,n)$ of $j$ and $k$ with respect to $l$.}
    \label{fig:AB_triangles}
\end{figure}


\section{Charge density due to applied electric potential}
\label{sec:micro-2nd}
When introducing a local electric potential, it breaks the bond inversion symmetry between $i,j$ for bonds with a potential difference. These bonds contribute to the charge density at second order.



The effective charge density operator is
\begin{equation}
        \delta \tilde{n}_l^{\phantom{(0)}}= \mathcal{P}_{1\text{h}} \, e^{S} \; \delta {n}_l \; e^{-S} \,  \mathcal{P}_{1\text{h}}
        = \delta n_l^{(0)} + \delta n_l^{(1)} + \delta n_l^{(2)} + \delta n_l^{(3)} + ...  \; ,
\end{equation}
where $\delta n_l^{(k)}$ is the contribution of $\delta \tilde{n}_l$ in $k$-th order of perturbation.
These contributions to the second order is
\begin{equation}
    \begin{split}
      %  \delta \tilde{n}_l^{\phantom{(0)}} &= \mathcal{P}_{1\text{h}} \, e^{S} \; \delta {n}_l \; e^{-S} \,  \mathcal{P}_{1\text{h}}
       % \\ &= \delta n_l^{(0)} + \delta n_l^{(1)} + \delta n_l^{(2)} + \delta n_l^{(3)} + ... \; , \\
        \delta n_l^{(0)} &= \mathcal{P}_{1\text{h}} \, \delta {n}_l \,\mathcal{P}_{1\text{h}} = 0  \; , \\
         \delta n_l^{(1)} &= \mathcal{P}_{1\text{h}} [S_1, \delta {n}_l ] \mathcal{P}_{1\text{h}} = 0\; , \\
         \delta n_l^{(2)} &=  %\frac{1}{2} \mathcal{P}_{1\text{h}} [S_1,[S_1, \delta {n}_l ]]  \mathcal{P}_{1\text{h}} = 0\; , \\
         \mathcal{P}_{1\text{h}} \Big\{\, [S_2, \delta {n}_l ]+  \frac{1}{2} [S_1,[S_1, \delta {n}_l ]] \; \Big\}  \mathcal{P}_{1\text{h}}  \; ,
    \end{split}
\end{equation}
where $S_1$ and $S_2$ are the generators of the canonical transformation for ${H}$ with
After projecting into the low-energy subspace, only the operators with zero total ``charge'' remain, meaning terms that do not change the hole occupancy. For example, $ [S_1,[S_1, [S_1, \delta {n}_l ]]] = 0$ since each $S_1$ changes the occupancy by $\pm 1$ and there is no way to combine three of these terms to keep within the same occupancy subspace.

The generator of the transformation $S=\sum_k S_k$ is written as $S_k=S_k^+-S_k^-$ with $S_k^{\pm}$ being a homogeneous polynomial of degree $k$ in $T_{1},T_0,T_{-1}$ that has charge $\pm 1$ with property $\mathcal{P}_{1\text{h}}S_k^+ = S_k^-\mathcal{P}_{1\text{h}} =0$. From this it follows directly that $\mathcal{P}_{1\text{h}} [S_k, \delta {n}_l ] \mathcal{P}_{1\text{h}} = 0$ and then $\delta n_l^{(0)}=\delta n_l^{(1)}=0$.
The second-order term can be expanded to
    \begin{equation}
\begin{split}
    \delta n_l^{(2)}
    \!&=\!  - \mathcal{P}_{1\text{h}} \, S_1^{-} \, [ \, {n}_l, S_1^{+} ]\,  \mathcal{P}_{1\text{h}} = \sum_{j} \delta n_{l \,(j)}^{(2)} \; ,
\end{split}
\end{equation} where %the sum is over NN of $l$ and
$\delta n_{l \,(j)}^{(2)}$  is the charge transfer from $j\to l$ minus the charge transfer from $l \to j$.
Including $V_0$ the charge density has a non-zero second order contribution.
The contribution for $\delta n_{l \,(j)}^{(2)}$ of the generator of the canonical transformation in the presence of the electric potential is
\begin{align}
    S_1^{+}  &\overset{\;i \,(j)\;}{\longrightarrow}  \sum_{\ell=0} \Big[ \frac{1}{\Delta E_{\ell}+eV_0} \mathcal{P}^{(2)}_{i,\ell} \, h_i^\dagger  \tilde{\bf{T}}_{ij} h_j^{\phantom{\dagger}} \, \mathcal{P}^{(1)}_{j}  + \frac{1}{\Delta E_{\ell}-eV_0} \mathcal{P}^{(2)}_{j,\ell} \, h_j^\dagger  \tilde{\bf{T}}_{ji} h_i^{\phantom{\dagger}} \, \mathcal{P}^{(1)}_{i} \Big] \; .
\end{align}
The second order change in  charge density
\begin{align}
     \delta n_{\;i \,(j)\;}^{(2)}   &= \sum_{\ell=0}^{2}\!\! \Big[ \frac{1}{(\Delta E_{\ell}+eV_0)^2}  h_j^\dagger  \tilde{\bf{T}}_{ji} h_i^{\phantom{\dagger}}\, \mathcal{P}^{(2)}_{i,\ell} \, h_i^\dagger  \tilde{\bf{T}}_{ij} h_j^{\phantom{\dagger}}- \frac{1}{(\Delta E_{\ell}-eV_0)^2} h_i^\dagger  \tilde{\bf{T}}_{ij} h_j^{\phantom{\dagger}} \mathcal{P}^{(2)}_{j,\ell} \, h_j^\dagger  \tilde{\bf{T}}_{ji} h_i^{\phantom{\dagger}}   \Big] \; , \\
     &= \sum_{\ell}\! \sum_{s_1, s_2, s_3,s_4}\!\! \Big[ \frac{ (\tilde{\bf{T}}_{ji})_{s_1,s_2}(\tilde{\bf{T}}_{ij})_{s_3,s_4}  }{(\Delta E_{\ell}+eV_0)^2}    h_{i,s_2}^{\phantom{\dagger}}\, \mathcal{P}^{(2)}_{i,\ell} \, h_{i,s_3}^\dagger
 h_{j,s_1}^\dagger h_{j,s_4}^{\phantom{\dagger}}- \big( i\leftrightarrow j,\, V_0\to-V_0\big) \Big] \; .
\end{align}
From this we confirm that $ \delta n_{\;i \,(j)\;}^{(2)}=- \delta n_{\;j \,(i)\;}^{(2)}  $.
The matrix elements are
\begin{align}
     \braket{ s_i',\,s_j'\,| \delta n_{\;i \,(j)\;}^{(2)} | \, s_i, \, s_j}   &=
      \sum_{\ell}\! \sum_{s_2, s_3}\! \Big[ \frac{ (\tilde{\bf{T}}_{ji})_{s_j',s_2}(\tilde{\bf{T}}_{ij})_{s_3,s_j}  }{(\Delta E_{\ell}+eV_0)^2}    F_{\ell}(s_i',s_2,s_3,s_i)
- \big( i\leftrightarrow j,\, V_0\to-V_0\big) \Big] \; .
\end{align}
where, as before, we have  $
    F_{\ell}(s_1,s_2,s_3,s_4) = \braket{s_1|h_{s_2
}^{\phantom{\dagger}}\, \mathcal{P}^{(2)}_{\ell} \, h_{s_3}^\dagger|s_4}$.



Including $V_0$ the charge density has a non-zero second order contribution. For a pair of sites in a $\gamma$ bond
\begin{equation}
    \delta n^{(2)}_{l\,(k)}
    =F_{\mathbb{1}}
    - \Bigg[   F_{J} \bm{\sigma}_l \cdot \bm{\sigma}_k
    + F_{K} \sigma_l^\gamma \sigma_k^\gamma
    + F_{\Gamma} \left( \sigma^\alpha_l\sigma^\beta_k + \sigma^\alpha_l\sigma^\alpha_k \right) \Bigg] \; ,
\end{equation}
where the $F$ are functions of $t_i$, $J_H$, $U$ and $e V_0$. The sub-indexes indicate a relation to the couplings $J, K, \Gamma$, in particular, the $t_i$ dependence are the same that appears in these coefficients. The functions $F_K$ and $F_\Gamma$ are multiples of $K$ and $\Gamma$, $F_K = K(V_0) \; g(\xi,\eta), F_\Gamma = \Gamma(V_0) \;  g(\xi,\eta)$, but $F_J$ and $F_\mathbb{1}$ are more complicated.


\begin{figure}[h!]
    \centering
    \includegraphics[width=0.7\textwidth]{Chapters/Appendix/Figures/micro/plot-second-order-t1-xi1-g0-hc.2.pdf}\\
    \includegraphics[width=0.7\textwidth]{Chapters/Appendix/Figures/micro/plot-second-order-t1-xi1-g4-hc.2.pdf}\\
    \includegraphics[width=0.7\textwidth]{Chapters/Appendix/Figures/micro/plot-third-order-t1-xi1-g4-hc.2.pdf}
    \caption{[Top and middle] \textbf{Second} + \textbf{Third} order charge density in presence of an electric potential $(\xi=0.1)$ in configuration without ($g0$) and with vortices ($g4$). For comparison, in Fig. [Bottom] only \textbf{Third} order contribution to the charge density for a configuration with a vortex is shown.}
    \label{fig:micro-2nd}
\end{figure}


As expected, the charge in response to the electric potential is orders of magnitude larger than the charge due to a single vortex.
The contribution from the vortices is a third-order perturbation effect, and it is typically of order $\sim 10^{-6} $ in units of electron charge.
While the contribution due to an electric potential is a third-order perturbation effect, and for potentials of order $10^{-1} \, \mathrm{V}$ contributes to a charge of order $\sim 10^{-3}$. This is shown in Fig.~\ref{fig:micro-2nd}, where we use an ultra-local potential that yields a change in the charge density in the first $3$ rings from the center of the plaquette where the potential is applied.

We find that the applied potential completely dominates over the charge due to vortices. Moreover, in the second order, the charge contribution for a configuration with and without vortices is essentially the same, allowing one to compute its difference and extract the contribution in the third order solely due to the vortices.




